\section{Clifford+T is universal}
\label{sec:universal}

This section assembles the main theorem from the exact synthesis of
Section~\ref{sec:exact} and the $R_z$ approximation of
Section~\ref{sec:lemma12}.  Both nodes below and the main theorem are proved;
the axiom closure of the Lean main theorem is exactly
\texttt{[propext, Classical.choice, Quot.sound]}, so it depends on no
\texttt{sorry} and on no \texttt{native\_decide} certificate.

\begin{lemma}[Embedded $R_z$ approximation]
  \label{lem:rz_embedded_approx}
  \lean{TwoControl.Clifford.Universal.embedded_rz_approximation_by_clifford_t}
  \leanok
  Let $R$ be an embedded $R_z(\theta)$ gate in an $n$-qubit register and let
  $\epsilon > 0$.  Then there is a circuit over $\{C(X),H,T\}$ with
  $d(R, \operatorname{circuitMatrix}) < \epsilon$.
\end{lemma}

\begin{proof}
  \leanok
  \uses{thm:lemma12, def:embedded_one, def:hs_distance, def:ht_primitive}
  Apply Theorem~\ref{thm:lemma12} on the target wire to get a one-qubit
  $\{H,T\}$ circuit within $\epsilon$ of $R_z(\theta)$, then lift it along the
  placement that defines the embedding.  The step that makes this work is that
  the Hilbert--Schmidt distance is unchanged by the embedding: reindexing and
  padding with identity factors multiply the relevant trace by
  $\operatorname{Tr}(I) = 2^{n-1}$ and the dimension by the same factor, so the
  normalized quantity is invariant.  There are three cases, according to
  whether the gate sits on a one-wire placement or as either half of a two-wire
  placement.
\end{proof}

\begin{lemma}[Replacing every $R_z$ in a circuit]
  \label{lem:rz_replacement}
  \lean{TwoControl.Clifford.Universal.clifford_rz_synthesis_approximates_by_clifford_t}
  \leanok
  Suppose $U$ is synthesized up to global phase by a circuit over
  $\{C(X),H,T,R_z(\theta)\}$.  Then for every $\epsilon > 0$ there is a circuit
  $C$ over $\{C(X),H,T\}$ with $d(U, C) < \epsilon$.
\end{lemma}

\begin{proof}
  \leanok
  \uses{lem:rz_embedded_approx, lem:hs_mul, lem:hs_phase,
        lem:hs_self, def:gate_cliffordtrz}
  Let the given circuit have $k$ factors and put $\delta = \epsilon/(k+1)$.
  Replace each factor by a circuit over $\{C(X),H,T\}$ within $\delta$: an
  exact Clifford+T factor is kept as is and contributes distance $0$ by
  Lemma~\ref{lem:hs_self}, and an embedded $R_z$ factor is replaced using
  Lemma~\ref{lem:rz_embedded_approx}.  Composing the replacements factor by
  factor and applying Lemma~\ref{lem:hs_mul} at each step bounds the total
  distance by $k\delta < \epsilon$.  Finally the global phase relating $U$ to
  the original circuit matrix contributes nothing, by
  Lemma~\ref{lem:hs_phase}.
\end{proof}

\begin{theorem}[Clifford+T is universal]
  \label{thm:clifford_t_universal}
  \lean{TwoControl.Clifford.Universal.clifford_t_is_universal}
  \leanok
  Let $n \ge 0$, let $U$ be a $2^n \times 2^n$ unitary, and let $\epsilon > 0$.
  Then there is a circuit $C$ over $\{C(X), H, T\}$ with
  \[
    d(U, C) < \epsilon .
  \]
\end{theorem}

\begin{proof}
  \leanok
  \uses{thm:clifford_rz_universal, thm:clifford_rz_to_t_rz,
        lem:rz_replacement}
  For $n \ge 1$: Theorem~\ref{thm:clifford_rz_universal} synthesizes $U$
  exactly over $\{C(X),H,S,S^\dagger,R_z\}$ up to global phase,
  Theorem~\ref{thm:clifford_rz_to_t_rz} rewrites $S = T^2$ and
  $S^\dagger = T^6$, and Lemma~\ref{lem:rz_replacement} replaces every
  remaining $R_z$ by a Clifford+T approximation.

  For $n = 0$ the state space is one dimensional, every unitary is a scalar of
  modulus one, and the empty circuit is at distance $0$.
\end{proof}
